Un générateur idéal de force électromotrice $E=\qty{12}{\V}$ est utilisé pour
alimenter le <<pont de résistances>> représenté sur le schéma ci-dessous.
\begin{figure}[H]
    \centering
    \includegraphics[width=.7\textwidth]{585c8fb5-5d8c-4d0d-afc7-b49e8de01aef-5_348_789_1074_104}
\end{figure}

\begin{enumerate}
    \item Déterminer la résistance équivalente à l'ensemble et en déduire
          l'intensité $I$ du courant débité par le générateur.
    \item Calculer les intensités $I_{1}$ et $I_{2}$ des courants qui circulent
          respectivement dans les branches $ADC$ et $ABC$.
    \item En déduire la tension aux bornes de chaque conducteur.
    \item Calculer les puissances consommées et vérifier la conservation de
          l'énergie électrique pour ce circuit.
    \item Calculer la tension $U_{BD}$.
    \item Que se passe-t-il si on place un fil conducteur entre ces deux
          points~?
    \item Calculer la nouvelle résistance équivalente et vérifier qu'elle est
          bien égale à la précédente.
\end{enumerate}

\begin{donnees}
    $R_{1}=\qty{20}{\ohm}$~; $R_{2}=\qty{5.0}{\ohm}$~;
    $R_{3}=\qty{15}{\ohm}$~; $R_{4}=60 \Omega$.
\end{donnees}

